\documentclass[11pt]{article}
\usepackage[margin=1in]{geometry}
\usepackage{amsmath,amssymb,amsthm}
\usepackage{booktabs}
\usepackage{enumitem}
\usepackage{hyperref}
\usepackage{xcolor}
\usepackage{algorithm}
\usepackage{algorithmic}

\definecolor{maincolor}{RGB}{0,70,130}
\hypersetup{colorlinks=true,linkcolor=maincolor,urlcolor=maincolor}

\newcommand{\bigO}{\mathcal{O}}
\newcommand{\R}{\mathbb{R}}
\newcommand{\CE}{\mathrm{CE}}
\DeclareMathOperator*{\argmin}{argmin}
\DeclareMathOperator*{\argmax}{argmax}

\newtheorem{definition}{Definition}
\newtheorem{proposition}{Proposition}
\newtheorem{remark}{Remark}

\title{\textbf{Compact Cutting-Plane for Congestion Games} \\ \large Implementation Specification}
\author{}
\date{}

\begin{document}
\maketitle

%==========================================================================
\section{Congestion Game Model}
%==========================================================================

\subsection{Network and players}

Consider a directed graph $G = (V, E)$ with $|E| = K$ edges. There are $n$ agents. Each agent $i$ has a source $s_i \in V$ and a destination $t_i \in V$. Agent $i$'s action set $A_i$ is a set of $s_i$--$t_i$ paths. We write $m_i = |A_i|$.

Each action $a_i \in A_i$ is a path, represented as a subset of edges $a_i \subseteq E$.

\subsection{Edge cost model}

Each edge $e \in E$ has a latency function parameterized by $\theta_e \in \R^p$. The latency on edge $e$ depends on the congestion $\ell_e(a)$ (number of agents using edge $e$ under action profile $a$):
\[
c_e\bigl(\ell_e(a)\bigr) \;=\; \sum_{j=1}^{p} \theta_{e,j} \cdot \ell_e(a)^{j-1} \;=\; \theta_e^\top \psi\bigl(\ell_e(a)\bigr)
\]
where $\psi(\ell) = (1, \ell, \ell^2, \ldots, \ell^{p-1})^\top \in \R^p$ is the polynomial basis.

\begin{remark}
For affine latencies ($p=2$): $c_e(\ell) = \theta_{e,1} + \theta_{e,2}\,\ell$, so $\theta_e = (\theta_{e,1}, \theta_{e,2})$.
\end{remark}

\subsection{Utility function}

Agent $i$'s utility at action profile $a = (a_1, \ldots, a_n)$ is the negative total latency on their path:
\[
u_i(a) \;=\; -\sum_{e \in a_i} c_e\bigl(\ell_e(a)\bigr)
\;=\; -\sum_{e \in a_i} \theta_e^\top \psi\bigl(\ell_e(a)\bigr)
\]

\subsection{Parameter vector}

Stack all edge parameters into a single vector:
\[
\theta \;=\; (\theta_1^\top, \theta_2^\top, \ldots, \theta_K^\top)^\top \;\in\; \R^d, \qquad d = Kp
\]

The utility is linear in $\theta$:
\[
u_i(a) \;=\; -\phi_i(a)^\top \theta
\]
where $\phi_i(a) \in \R^d$ is the \textbf{feature vector} for agent $i$ at profile $a$:
\[
\phi_i(a) \;=\; \bigl(\underbrace{\mathbf{1}[e \in a_i] \cdot \psi(\ell_e(a))}_{\text{block for edge } e}\bigr)_{e \in E} \;\in\; \R^d
\]
Blocks corresponding to edges \emph{not} in $a_i$'s path are zero.

\subsection{Utility differences}

The utility difference (the object the cutting-plane learns) is:
\[
w_i(a_i, a_i', a_{-i}) \;=\; u_i(a_i', a_{-i}) - u_i(a_i, a_{-i}) \;=\; -\bigl(\phi_i(a_i', a_{-i}) - \phi_i(a_i, a_{-i})\bigr)^\top \theta
\]
Define the \textbf{difference feature}:
\[
\Delta\phi_i(a_i, a_i', a_{-i}) \;:=\; \phi_i(a_i, a_{-i}) - \phi_i(a_i', a_{-i}) \;\in\; \R^d
\]
so that $w_i(a_i, a_i', a_{-i}) = \Delta\phi_i(a_i, a_i', a_{-i})^\top \theta$.

In the full normal-form representation, the packed $w$-vector satisfies $w = \Phi\,\theta$ where $\Phi \in \R^{N \times d}$ stacks all $\Delta\phi_i$ vectors. But we \textbf{never form $\Phi$ explicitly}.

%==========================================================================
\section{Compact Correlated Equilibrium LP}
%==========================================================================

\subsection{Why the standard CE LP is too large}

The standard CE LP has $M = \prod_i m_i$ variables (one per action profile). For $n=20$ agents with $m_i = 3$ routes each, $M = 3^{20} \approx 3.5 \times 10^9$. Intractable.

\subsection{Route-flow formulation}

Instead of variables $x(a)$ for each profile $a \in A$, we use a \textbf{route-flow} formulation. Define:
\[
f_{i,r} \;=\; \Pr_{a \sim x}[a_i = r] \qquad \text{for each agent } i, \text{ route } r \in A_i
\]
This is the marginal probability that agent $i$ plays route $r$. The number of variables is $\sum_i m_i$ (polynomial in $n$).

\begin{remark}
The marginals $f_{i,r}$ do not determine the joint distribution $x$ uniquely. However, for congestion games, the CE constraints depend on $x$ only through certain \textbf{congestion moments} (expected congestion on edges), which can be expressed in terms of the marginals under a product or low-correlation approximation. For \textbf{exact} CE computation in congestion games, one can use the approach of Papadimitriou and Roughgarden (2008), which exploits the structure of congestion games to solve the CE LP in polynomial time.
\end{remark}

\subsection{CE constraints in the congestion game}

For agent $i$, recommended route $r \in A_i$, deviation route $r' \in A_i$, the CE incentive constraint is:
\[
\sum_{a_{-i}} x(r, a_{-i}) \bigl[u_i(r', a_{-i}) - u_i(r, a_{-i})\bigr] \;\leq\; 0
\]
Substituting the linear utility model:
\[
\sum_{a_{-i}} x(r, a_{-i})\; \Delta\phi_i(r, r', a_{-i})^\top \theta \;\leq\; 0
\]

This can be written as:
\[
\Bigl(\sum_{a_{-i}} x(r, a_{-i})\; \Delta\phi_i(r, r', a_{-i})\Bigr)^\top \theta \;\leq\; 0
\]

The vector in parentheses is the \textbf{expected difference feature} under the conditional distribution $x(\cdot \mid a_i = r)$.

\subsection{Product-distribution CE approximation}

When the joint distribution $x$ is a product of marginals $x(a) = \prod_i f_{i, a_i}$, the expected congestion on edge $e$ is:
\[
\bar{\ell}_e \;=\; \sum_{i=1}^n \sum_{r \in A_i} f_{i,r}\, \mathbf{1}[e \in r]
\]
and the expected difference feature can be computed from $\bar{\ell}_e$ and the marginals without enumerating all profiles.

\textbf{Compact LP variables:} $f_{i,r}$ for $i \in [n]$, $r \in A_i$. Total: $\sum_i m_i$.

\textbf{Compact LP constraints:}
\begin{itemize}[nosep]
    \item Simplex per agent: $\sum_{r \in A_i} f_{i,r} = 1$, $f_{i,r} \geq 0$.
    \item CE incentive: for each $i$, $r$, $r'$: expected gain from deviating $\leq 0$. Under product distribution, these involve $\bar{\ell}_e$ and are polynomial-size.
\end{itemize}

\begin{remark}
The product-distribution approximation yields a \textbf{coarse correlated equilibrium} (CCE), not an exact CE. For exact CE, one needs the Papadimitriou-Roughgarden approach. For the experiments, CCE is a reasonable starting point and is much simpler to implement.
\end{remark}

%==========================================================================
\section{Compact Cutting-Plane Pipeline}
%==========================================================================

The full pipeline operates in $\R^d$ (the parameter space) instead of $\R^N$ (the normal-form utility space).

\subsection{State}

\begin{itemize}[nosep]
    \item $\theta(t) \in \R^d$: current estimate of edge cost parameters.
    \item Knowledge set $C_t = \{\theta \in [-R, R]^d : \langle q_\theta^{(s)}, \theta \rangle \geq 0,\; s = 1, \ldots, t\}$.
    \item Accumulated cuts $q_\theta^{(1)}, \ldots, q_\theta^{(t)} \in \R^d$.
\end{itemize}

\subsection{Per-round operations}

\textbf{Step 1: Query point.} Compute $\theta(t) = $ analytic center of $C_{t-1}$ via ACCPM in $\R^d$.

Cost: $\bigO(d \cdot t^2)$ per Newton step via Woodbury.

\medskip
\textbf{Step 2: Reconstruct utilities.} From $\theta(t)$, compute the utility differences needed for the CE LP. For each agent $i$, route pair $(r, r')$, and edge $e$:
\[
w_i(r, r')_e \;=\; \Delta\phi_i(r, r', \cdot)^\top \theta(t)
\]
This requires evaluating $\Delta\phi_i$ for each $(i, r, r', e)$ but \textbf{not} for each opponent profile $a_{-i}$ separately.

Under the product-distribution CE, the expected congestion depends on the marginals, which are LP variables. So we embed $\theta(t)$ into the CE LP constraints directly.

Cost: $\bigO(n \cdot m_{\max}^2 \cdot K)$ where $m_{\max} = \max_i m_i$.

\medskip
\textbf{Step 3: Solve compact CE LP.} Find $f = (f_{i,r})$ that forms a (coarse) correlated equilibrium of the game with parameters $\theta(t)$.

Number of variables: $\sum_i m_i$.\\
Number of constraints: $\sum_i m_i(m_i - 1)$ incentive constraints + $n$ simplex constraints.

Cost: Gurobi LP in $\sum_i m_i$ variables. Very fast.

\medskip
\textbf{Step 4: Sample and recommend.} Sample a route $r_i \sim f_{i,\cdot}$ independently for each agent $i$. Send $r_i$ privately to agent $i$.

Cost: $\bigO(n)$.

\medskip
\textbf{Step 5: Observe feedback.} Each agent $i$, receiving recommendation $r_i$ and knowing $f$, best-responds (or quantal-responds). The moderator observes the realized routes $r_i^*$ for all agents.

The realized action profile $a^* = (r_1^*, \ldots, r_n^*)$ determines the congestion $\ell_e(a^*)$ on each edge.

\medskip
\textbf{Step 6: Construct cut in $\R^d$.} For each deviating agent $i$ ($r_i^* \neq r_i$):

The separation oracle (Algorithm~4 of the paper) produces a cut $q \in \R^N$. But we compute the projected cut directly in $\R^d$ without ever forming $q$:
\[
q_\theta \;=\; \sum_{i:\, r_i^* \neq r_i} x(r_i, \cdot)^\top \Delta\Phi_i(r_i, r_i^*)
\]
where $x(r_i, \cdot)$ is the conditional distribution over opponents given $r_i$, and $\Delta\Phi_i(r_i, r_i^*)$ stacks the difference features.

Under the product distribution, $x(r_i, a_{-i}) = \prod_{j \neq i} f_{j, a_j}$, so:
\[
q_\theta \;=\; \sum_{i:\, r_i^* \neq r_i} \sum_{a_{-i}} \Bigl(\prod_{j \neq i} f_{j, a_j}\Bigr) \Delta\phi_i(r_i, r_i^*, a_{-i})
\]

\textbf{Key simplification for congestion games.} The feature $\Delta\phi_i(r_i, r_i^*, a_{-i})$ depends on $a_{-i}$ only through the congestion $\ell_e(a_{-i})$ on edges in $r_i \cup r_i^*$. The expected congestion under the product distribution is:
\[
\mathbb{E}[\ell_e(a_{-i})] \;=\; \sum_{j \neq i} \sum_{r \in A_j} f_{j,r}\, \mathbf{1}[e \in r]
\]
For the polynomial basis $\psi(\ell) = (1, \ell, \ldots, \ell^{p-1})$, we need the moments $\mathbb{E}[\ell_e^k]$ for $k = 0, \ldots, p-1$. Under a product distribution, these moments can be computed in $\bigO(n)$ per edge via the Poisson binomial distribution.

For \textbf{affine latencies} ($p=2$), we only need $\mathbb{E}[\ell_e]$ and $\mathbb{E}[\ell_e^2]$:
\begin{align*}
\mathbb{E}[\ell_e] &= \sum_{j \neq i} p_{j,e}, \qquad p_{j,e} := \sum_{r \in A_j} f_{j,r}\,\mathbf{1}[e \in r] \\
\mathbb{E}[\ell_e^2] &= \mathbb{E}[\ell_e]^2 + \sum_{j \neq i} p_{j,e}(1 - p_{j,e})
\end{align*}

This gives us $q_\theta \in \R^d$ directly, without enumerating opponent profiles.

Cost: $\bigO(n \cdot K \cdot p)$.

\medskip
\textbf{Step 7: Update knowledge set.} Add $q_\theta$ to the cut list. Compute next analytic center.

\subsection{Separation property}

The projected cut preserves the separation property:
\begin{align*}
\langle \theta^*, q_\theta \rangle &= \langle \Phi\,\theta^*, q \rangle_{\R^N} = \sum_i \phi_i(r_i, r_i^*, x) \;\geq\; 0 \\
\langle \theta(t), q_\theta \rangle &= \langle \Phi\,\theta(t), q \rangle_{\R^N} \;\leq\; 0
\end{align*}
where the first inequality holds because agents deviate to profitable actions, and the second because $x(t) \in \CE(\theta(t))$.

%==========================================================================
\section{Complexity Summary}
%==========================================================================

\begin{table}[h]
\centering
\begin{tabular}{lcc}
\toprule
\textbf{Operation} & \textbf{Normal-form} & \textbf{Congestion (compact)} \\
\midrule
ACCPM (per Newton) & $\bigO(N \cdot t^2)$ & $\bigO(d \cdot t^2)$ \\
CE LP variables & $M = \prod_i m_i$ & $\sum_i m_i$ \\
CE LP constraints & $\sum_i m_i(m_i-1)$ & $\sum_i m_i(m_i-1)$ \\
Cut construction & $\bigO(M)$ & $\bigO(n \cdot K \cdot p)$ \\
\midrule
$N$ (parameter dim) & $\sum_i(m_i{-}1)|A_{-i}|$ & $d = Kp$ \\
\bottomrule
\end{tabular}
\caption{Per-round cost comparison. $K$ = edges, $p$ = polynomial degree, $d = Kp$, $t$ = accumulated cuts, $M = \prod_i m_i$, $N = $ normal-form utility dimension.}
\end{table}

\textbf{Example:} $n = 100$ agents on a network with $K = 50$ edges, affine latencies ($p=2$), $m_i = 5$ routes each.
\begin{itemize}[nosep]
    \item Normal-form: $M = 5^{100} \approx 10^{70}$. \textbf{Completely intractable.}
    \item Compact: $d = 100$, CE LP has 500 variables. \textbf{Trivial.}
\end{itemize}

%==========================================================================
\section{Implementation Plan}
%==========================================================================

\subsection{Module: \texttt{congestion\_game.py}}

\begin{enumerate}[nosep]
    \item \textbf{Network definition:} Graph $G = (V, E)$, source-destination pairs $(s_i, t_i)$, route sets $A_i$ (enumerate $s_i$--$t_i$ paths up to a length bound).
    \item \textbf{Feature computation:} Given an action profile $a$ and agent $i$, compute $\phi_i(a) \in \R^d$ and $\Delta\phi_i(r, r', a_{-i}) \in \R^d$.
    \item \textbf{Utility evaluation:} $u_i(a) = -\phi_i(a)^\top \theta$.
    \item \textbf{Congestion computation:} Given marginals $f$, compute expected congestion moments $\mathbb{E}[\ell_e^k]$ per edge.
\end{enumerate}

\subsection{Module: \texttt{compact\_ce\_solver.py}}

\begin{enumerate}[nosep]
    \item \textbf{Build LP:} Variables $f_{i,r} \geq 0$, simplex constraints, CE incentive constraints using $\theta(t)$ and expected congestion moments.
    \item \textbf{Solve via Gurobi:} LP with $\sum_i m_i$ variables.
    \item \textbf{Return marginals:} $f_{i,r}$ for each agent $i$ and route $r$.
\end{enumerate}

\subsection{Module: \texttt{compact\_oracle.py}}

\begin{enumerate}[nosep]
    \item \textbf{Sample:} Route $r_i \sim f_{i,\cdot}$ independently per agent.
    \item \textbf{Feedback:} Each agent best-responds (or quantal-responds) to the recommendation given the marginals. Observe realized routes $r_i^*$.
    \item \textbf{Cut in $\R^d$:} For each deviating agent, compute $q_\theta$ using expected congestion moments. No enumeration of $a_{-i}$.
\end{enumerate}

\subsection{Module: \texttt{compact\_cutting\_plane.py}}

\begin{enumerate}[nosep]
    \item \textbf{Reuse existing ACCPM} from \texttt{cutting\_plane.py}, but in $\R^d$ instead of $\R^N$.
    \item \textbf{Per-round loop:} $\theta(t)$ from ACCPM $\to$ compact CE LP $\to$ sample $\to$ feedback $\to$ cut in $\R^d$ $\to$ update.
    \item \textbf{Track regret:} Both platform compliance regret and player swap regret, same as before.
\end{enumerate}

\subsection{Experiment}

Run the compact pipeline on a Braess-like network:
\begin{itemize}[nosep]
    \item $n = 10, 50, 100, 500$ agents.
    \item $K = 10$--$50$ edges, affine latencies ($p=2$, so $d = 2K$).
    \item $m_i = 3$--$5$ routes per agent.
    \item Compare with the normal-form pipeline (for small $n$ where both are feasible).
    \item Show that regret scaling matches theory and that the compact pipeline handles $n = 500$ in seconds.
\end{itemize}

\end{document}
